% from ega4-18-fr.tex
\oldpage[IV-4]{111}

\begin{env}[18.2.1]
\EGAReaderTarget{ega4.c04.statement.18-2-1}{IV.18.2.1}
Étant donnés un anneau $A$ et une $A$-algèbre commutative $B$ qui est \emph{finie} et est un $A$-module \emph{libre}, rappelons (Bourbaki, \emph{Alg.}, chap. VIII, \textsection~12, n\textsuperscript{o}~2) que l'on définit sur $B$ une \emph{forme $A$-linéaire} $\operatorname{Tr}_{B/A}$, la « forme trace »; on en déduit la définition d'une \emph{forme $A$-bilinéaire symétrique} (dite encore « forme trace »)
\begin{equation}
  (x,y)\longmapsto\operatorname{Tr}_{B/A}(xy)
  \tag{18.2.1.1}\label{IV.18.2.1.1}
\end{equation}
dont la donnée est équivalente à celle de l'application $A$-linéaire associée $\operatorname{astr}_{B/A}:B\to\check B$ du $A$-module $B$ dans son \emph{dual} $\check B$, égale à sa transposée. Lorsque $A$ est un \emph{corps}, il est équivalent de dire que cette forme bilinéaire est \emph{non dégénérée} ou que $B$ est une $A$-algèbre \emph{séparable} (Bourbaki, \emph{Alg.}, chap. IX, \textsection~2, prop. 5).

Soit $f:A\to A'$ un homomorphisme d'anneaux; posons $B'=B\otimes_A A'$, et soit $g:B\to B'$ l'homomorphisme canonique; l'image par $g$ d'une base du $A$-module $B$ est alors une base du $A'$-module $B'$, et il résulte des définitions que l'on a, pour tout $x\in B$,
\begin{equation}
  \operatorname{Tr}_{B'/A'}(g(x))=f(\operatorname{Tr}_{B/A}(x)).
  \tag{18.2.1.2}\label{IV.18.2.1.2}
\end{equation}
\end{env}

\begin{env}[18.2.2]
\EGAReaderTarget{ega4.c04.statement.18-2-2}{IV.18.2.2}
Considérons maintenant un \emph{espace annelé} $(X,\sh{O}_X)$ et soit $\sh{B}$ une $\sh{O}_X$-Algèbre qui, en tant que $\sh{O}_X$-Module, est \emph{localement libre de rang fini}; alors, pour tout ouvert $U\subset X$ tel que $\sh{B}|U$ soit (en tant que $\sh{O}_U$-Module) isomorphe à $\sh{O}_U^n$ (pour un $n$ dépendant de $U$), $\Gamma(U,\sh{B})$ est une $\Gamma(U,\sh{O}_X)$-algèbre qui, en tant que $\Gamma(U,\sh{O}_X)$-module, est libre de rang fini et définit donc une forme $\Gamma(U,\sh{O}_X)$-linéaire $\operatorname{Tr}_{\Gamma(U,\sh{B})/\Gamma(U,\sh{O}_X)}$, que nous noterons aussi $\operatorname{Tr}_{\sh{B}|\sh{O}_X,U}$; on en déduit une application linéaire associée
\[
  \operatorname{astr}_{\sh{B}|\sh{O}_X,U}:\Gamma(U,\sh{B})\longrightarrow\Gamma(U,\sh{B})^\vee=\Gamma(U,\check{\sh{B}}).
\]
En outre, il résulte de (18.2.1.2) que ces applications linéaires sont compatibles avec les opérations de restriction de $U$ à un ouvert plus petit, et définissent donc, d'une part un \emph{homomorphisme de $\sh{O}_X$-Modules}, appelé encore \emph{homomorphisme trace}:
\begin{equation}
  \operatorname{Tr}_{\sh{B}|\sh{O}_X}:\sh{B}\longrightarrow\sh{O}_X
  \tag{18.2.2.1}\label{IV.18.2.2.1}
\end{equation}
et d'autre part un \emph{homomorphisme de $\sh{O}_X$-Modules}
\begin{equation}
  \operatorname{astr}_{\sh{B}|\sh{O}_X}:\sh{B}\longrightarrow\check{\sh{B}}
  \tag{18.2.2.2}\label{IV.18.2.2.2}
\end{equation}
dit \emph{associé à la trace}, et égal à son transposé. Il résulte aussi de (18.2.1.2) que pour tout $x\in X$, on a
\begin{equation}
  (\operatorname{Tr}_{\sh{B}|\sh{O}_X})_x=\operatorname{Tr}_{\sh{B}_x|\sh{O}_{X,x}}
  \tag{18.2.2.3}\label{IV.18.2.2.3}
\end{equation}
\begin{equation}
  (\operatorname{astr}_{\sh{B}|\sh{O}_X})_x=\operatorname{astr}_{\sh{B}_x|\sh{O}_{X,x}}.
  \tag{18.2.2.4}\label{IV.18.2.2.4}
\end{equation}
Enfin, sous les conditions de (18.2.1), si l'on pose $X=\Spec(A)$ et si $\sh{B}=\widetilde B$ est la $\sh{O}_X$-Algèbre correspondant à $B$, la forme $\operatorname{Tr}_{\sh{B}|\sh{O}_X}$ (resp. l'homomorphisme $\operatorname{astr}_{\sh{B}|\sh{O}_X}$) correspond à la forme $\operatorname{Tr}_{B/A}$ (resp. à l'homomorphisme de $A$-modules $\operatorname{astr}_{B/A}$), comme il résulte encore de (18.2.1.2).
\end{env}


